%% $Id: pst-fit-doc.tex 881 2024-03-17 20:32:53Z herbert $
\documentclass[11pt,english,BCOR10mm,DIV12,bibliography=totoc,parskip=false,
   smallheadings, headexclude,footexclude,oneside]{pst-doc}
\usepackage[utf8]{inputenc}
\usepackage{pst-fit}
\let\pstFitFV\fileversion

\usepackage{biblatex}
\addbibresource{\jobname.bib}

\input{pst-fit-doc.data}
\readdata[ignoreLines=0]{\Gauss}{gauss.dat}
\readdata[ignoreLines=0]{\Power}{power.dat}
\readdata[ignoreLines=0]{\King}{king.dat}
\readdata[ignoreLines=0]{\Exp}{exp.dat}
\readdata[ignoreLines=0]{\Linear}{linear.dat}
\readdata[ignoreLines=0]{\LinearA}{linearA.dat}
\readdata[ignoreLines=0]{\Recip}{recip.dat}
\readdata[ignoreLines=0]{\Poly}{poly.dat}
\readdata[ignoreLines=0]{\Logt}{logt.dat}
\readdata[ignoreLines=0]{\Loge}{loge.dat}
\readdata[ignoreLines=0]{\Calib}{calib.dat}
\readdata[ignoreLines=0]{\PScatter}{pscatter.dat}
\readdata[ignoreLines=0]{\HighOrder}{highorder.dat}


\renewcommand\bgImage{%
\begin{psgraph}[arrows=<->,Dy=20](0,0)(0,0)(10,100){2in}{2in}
\listplot[decimals=2,EqPos=5 50,plotstyle=GLLSR,RecipFit,linestyle=dashed]{\Recip}
\listplot[plotstyle=dots]{\Recip}
\end{psgraph}
}

\let\Section\section
\def\section{\clearpage\Section}
% prints the string which the PostScript code leaves on the stack,
% the code is run with the procedures of the bignum library
\newcommand\bnshow[1]{%
  \begin{pspicture}(0,-0.1)(11,0.3)
  \pscustom{\moveto(0,0)
    \code{/Courier findfont 9 scalefont setfont tx@NumericDict begin #1 end show}}
  \end{pspicture}}
\newcommand\wrong[1]{\textcolor{gray}{#1}}% wrong digits
\begin{document}

\title{\texttt{pst-fit}}
\subtitle{Curve fitting; v.\pstFitFV}
\author{Buddy Ledger\\
Herbert Voß}
\docauthor{}
\date{\today}
\maketitle

\tableofcontents
%
%
%
\psset{xAxisLabel=x, %Setup for psgraph env
        yAxisLabel=y,
        xAxisLabelPos=,
        yAxisLabelPos=,
        lly=-1cm,llx=-1cm}

\section{Fit of Linear Function}
%Replaces LSM using default options

\begin{LTXexample}[pos=t]
\begin{psgraph}[arrows=->,Dy=2](0,0)(0,0)(10,22){4.5in}{4.5in}
\listplot[decimals=2,EqPos=3 15,plotstyle=GLLSR,linestyle=dashed]{\Linear}
\listplot[plotstyle=dots,linecolor=red,dotscale=2,dotstyle=x]{\Linear}
\end{psgraph}
\end{LTXexample}

\section{Fit of Power Function}
\begin{LTXexample}[pos=t]
\begin{psgraph}[arrows=->,Dy=100](0,0)(0,0)(10,1000){4.5in}{4.5in}
\listplot[decimals=2,EqPos=1 800,linestyle=dashed,plotstyle=GLLSR,PowerFit]{\Power}
\listplot[plotstyle=dots,linecolor=red,dotscale=2,dotstyle=x]{\Power}
\listplot[EqPos=5 4,plotstyle=GLLSR,linecolor=blue,CustomFit,CheckZeroX,CheckZeroY,FYtrans=log,RYtrans=10 exch exp,FXtrans=log]{\Power}
\end{psgraph}
\end{LTXexample}

\section{Fit of exp Function}
\begin{LTXexample}[pos=t]
\begin{psgraph}[arrows=->,Dx=1,Dy=1000,xsubticks=1,ysubticks=1](0,0)(0,0)(10,10000){4.5in}{4.5in}
\listplot[PstDebug=1,decimals=2,EqPos=4 2000,MaPos=5 7000,plotstyle=GLLSR,ExpFit]{\Exp}
\listplot[plotstyle=dots]{\Exp}
\end{psgraph}
\end{LTXexample}

\section{Fit of Log10/Loge Functions}
\begin{LTXexample}[pos=t]
\begin{psgraph}[arrows=->,Dx=1,Dy=2,xsubticks=1,ysubticks=1](0,0)(0,0)(10,50){4.5in}{4.5in}
\listplot[decimals=2,EqPos=5 10,plotstyle=GLLSR,LogTFit]{\Logt}
\listplot[decimals=2,plotstyle=dots]{\Logt}
\listplot[decimals=2,EqPos=5 28,plotstyle=GLLSR,LogEFit]{\Loge}
\listplot[decimals=2,plotstyle=dots]{\Loge}
%\listplot[EqPos=5 4,plotstyle=GLLSR,CustomFit,CheckZeroX=true,CheckZeroY=true,FYtrans=Yint dup mul sub %log,RYtrans=10 exch exp Yint dup mul add,FXtrans=log]{\Linear}
\end{psgraph}
\end{LTXexample}



\section{Fit of Recip}
\begin{LTXexample}[pos=t]
\begin{psgraph}[arrows=<->,Dx=1,Dy=10,xsubticks=1,ysubticks=1](0,0)(0,0)(10,100){4.5in}{4.5in}
\listplot[decimals=2,EqPos=5 2,plotstyle=GLLSR,RecipFit]{\Recip}
\listplot[plotstyle=dots]{\Recip}
\end{psgraph}
\end{LTXexample}

\section{Fit of Kings Law data}

\begin{LTXexample}[pos=t]
\begin{psgraph}[arrows=->,Dx=1,Dy=1,xsubticks=1,ysubticks=1](0,0)(0,0)(10,20){4.5in}{4.5in}
\pstScalePoints(1,1){}{dup mul}
\listplot[decimals=2,EqPos=5 2,plotstyle=GLLSR,PowerFit,Yint=5]{\King}
\listplot[plotstyle=dots]{\King}
\pstScalePoints(1,1){}{}
\end{psgraph}
\end{LTXexample}

\section{Fit of Guassian}
\begin{LTXexample}[pos=t]
\begin{psgraph}[arrows=->,Dx=1,Dy=0.1,xsubticks=1,ysubticks=1](0,0)(-10,0)(10,1){4.5in}{4.5in}
\listplot[plotstyle=dots]{\Gauss}
\listplot[decimals=2,EqPos=5 0.4,plotstyle=GLLSR,GaussFit,plotpoints=400]{\Gauss}
\end{psgraph}
\end{LTXexample}

\section{Fit of 4th Order Polynomial}
\begin{LTXexample}[pos=t]
\begin{psgraph}[arrows=->,Dx=1,Dy=100,xsubticks=1,ysubticks=1](0,0)(0,-600)(10,600){4.5in}{4.5in}
\listplot[plotstyle=dots]{\Poly}
%note that the valuewidth needs to be large for debugging matrices
\listplot[valuewidth=20,PstDebug=1,decimals=2,EqPos=0.5 -200,plotstyle=GLLSR,MaPos=7 500,MaScale=0.5,PolyOrder=4,plotpoints=400,Yint=500]{\Poly}
\listplot[linestyle=dashed,decimals=2,EqPos=0.5 -400,plotstyle=GLLSR,PolyOrder=4,plotpoints=400]{\Poly}
%\listplot[decimals=2,EqPos=0.5 -200,plotstyle=GLLSR,PolyOrder=10,plotpoints=400]{\Power}
\end{psgraph}
\end{LTXexample}

\section{LinetoXAxis respects scalepoints.}
\begin{LTXexample}[pos=t]
\begin{psgraph}[arrows=<->,Dx=1,Dy=1,xsubticks=1,ysubticks=1](0,0)(0,0)(10,10){4in}{4in}
\pstScalePoints(0.5,0.25){2 add}{16 add}
\listplot[plotstyle=line]{\Linear}
\listplot[arrows=<-,plotstyle=LineToXAxis]{\Linear}
\pstScalePoints(1,1){}{}
\rput[lb](1,2){Scaled fluid velocity distribution on a plot of the physical system.}
\end{psgraph}
\end{LTXexample}

\section{Prepare Points Modification}
\begin{LTXexample}[pos=t]
\begin{psgraph}[arrows=->,Dx=1,Dy=1,xsubticks=1,ysubticks=1](0,0)(0,0)(15,30){4.5in}{4.5in}
\listplot[plotstyle=line,plotNoMax=2,plotNo=1]{\LinearA}
\listplot[plotstyle=line,plotNoMax=2,plotNo=2]{\LinearA}
\listplot[plotstyle=line,linestyle=dashed,plotNoMax=2,plotNo=1,plotNoTwo=2,plotNoTwoFunc=add 2 div]{\LinearA}
\rput[lb](10,18){1st Column}
\rput[lb](10,29){2nd Column}
\rput[lb](10,24){Average of 1st and 2nd Columns}
\end{psgraph}
\end{LTXexample}

\section{PrintCoor Demo}
\begin{LTXexample}[pos=t]
\begin{psgraph}[arrows=->,Dx=1,Dy=1,xsubticks=1,ysubticks=1](0,0)(0,0)(10,15){4.5in}{4.5in}
%PrintCoor Demo
\psset{xEnd=,xStart=,decimals=2,valuewidth=50,fontscale=10,PstDebug=1}
\pstScalePoints(1,1){0.75 mul}{0.5 exp 3 mul}
\listplot[plotstyle=PrintCoor,plotNoMax=1,plotNo=1,PstDebug=0,relxerr=0,relyerr=0.1]{\Linear}
\listplot[plotstyle=dots,plotNoMax=1,plotNo=1]{\Linear}
\pstScalePoints(1,1){}{}
\end{psgraph}
\end{LTXexample}

\section{Regression statistics}
The least squares fit is linear in its coefficients; the method is described
in \cite{chapra98}. The statistical properties of the result, on which this
and the next two sections are based, are treated in \cite{dunn10},
\cite{devore99} and \cite{draper98}, in order of increasing rigour.
For the special fit
functions the data is first transformed, e.\,g.\ with the logarithm, and the
statistics below are those of the transformed data, where the fit is a straight
line or a polynomial (see \Lkeyword{plotcoord} for the alternatives).
With $n$ data points, $p$ fitted coefficients and the fitted values $\hat y_i$:
\begin{align*}
S_{xx} &= \sum (x_i-\bar x)^2 &
S_{yy} &= \sum (y_i-\bar y)^2 &
S_{xy} &= \sum (x_i-\bar x)(y_i-\bar y)\\
\mathit{SSE} &= \sum (y_i-\hat y_i)^2 &
R^2 &= 1-\frac{\mathit{SSE}}{S_{yy}} &
S_{yx} &= \sqrt{\frac{\mathit{SSE}}{S_{\mathrm{dof}}}},\quad S_{\mathrm{dof}}=n-p
\end{align*}
$S_{yx}$ is the standard error of the fit and $S_{\mathrm{dof}}$ the number of
degrees of freedom. The optional argument \Lkeyword{printconfig} takes up to 11 flags
(\texttt{0} or \texttt{1}) which select the lines to print below each other at
\Lkeyword{EqPos}. The default is \texttt{1 1 0 0 0 0 0 0 0 0 0}, the equation
and $R^2$.

\medskip
\begin{tabular}{@{}cl@{\qquad}cl@{}}\hline
flag & printed line & flag & printed line\\\hline
1 & equation of the fit                   & 7  & $S_{xy}$\\
2 & $R^2$                                 & 8  & $S_{yx}$\\
3 & Student $t$ value $t_{v,P}$           & 9  & mean value $\bar x$\\
4 & degrees of freedom $S_{\mathrm{dof}}$ & 10 & mean value $\bar y$\\
5 & $S_{xx}$                              & 11 & half width $\delta y_c$ of the confidence interval\\
6 & $S_{yy}$                              &    & at the first $x$, at $\bar x$ and at the last $x$\\\hline
\end{tabular}

\begin{LTXexample}[pos=t]
\begin{psgraph}[arrows=->,Dy=2](0,0)(0,0)(10,22){4.5in}{4.5in}
\listplot[decimals=3,EqPos=1 21,plotstyle=GLLSR,linestyle=dashed,
  printconfig=1 1 1 1 1 1 1 1 1 1 1]{\Calib}
\listplot[plotstyle=dots,linecolor=red,dotscale=2,dotstyle=x]{\Calib}
\end{psgraph}
\end{LTXexample}

\section{Confidence and prediction intervals}
The optional argument \Lkeyword{plottype} selects what the plot style \texttt{GLLSR} draws:

\medskip
\begin{tabular}{@{}ll@{}}\hline
\Lkeyword{plottype} & meaning\\\hline
\texttt{fit} & the fitted curve (default)\\
\texttt{upperci}, \texttt{lowerci}, \texttt{bothci} & limits of the confidence interval of the fit\\
\texttt{upperpi}, \texttt{lowerpi}, \texttt{bothpi} & limits of the prediction interval\\
\texttt{residuals} & the standardised residuals\\\hline
\end{tabular}

\medskip
The confidence interval of the fit (of the mean response) is the range in
which the fitted curve itself lies with a given probability. The prediction
interval is the range in which a single new measurement is expected, it is wider:
\[
\hat y(x) \pm t_{v,P}\,S_{yx}\sqrt{h(x)}
\qquad\text{and}\qquad
\hat y(x) \pm t_{v,P}\,S_{yx}\sqrt{1+h(x)}
\]
For a straight line the leverage is $h(x)=1/n+(x-\bar x)^2/S_{xx}$; for a
polynomial it is calculated from the inverse of the least squares matrix
\cite{draper98}.
$t_{v,P}$ is the two sided 95\,\% value of the Student $t$ distribution for
$S_{\mathrm{dof}}$ degrees of freedom, which is taken from a table
(see e.\,g.\ \cite{devore99}). Another
value can be set with the optional argument \Lkeyword{tvp}, e.\,g.\ \texttt{tvp=3.355}
for 99\,\% and 8 degrees of freedom. For the intervals use \texttt{ShowEq=false},
otherwise the equation is printed again.

\begin{LTXexample}[pos=t]
\begin{psgraph}[arrows=->,Dy=2](0,0)(0,0)(10,22){4.5in}{4.5in}
\psset{plotstyle=GLLSR,ShowEq=false}
\listplot[ShowEq,decimals=2,EqPos=1 20]{\Calib}
\listplot[plottype=bothci,linestyle=dashed,linecolor=blue]{\Calib}
\listplot[plottype=bothpi,linestyle=dotted,linecolor=red]{\Calib}
\listplot[plotstyle=dots]{\Calib}
\end{psgraph}
\end{LTXexample}

The intervals are also available for the other fit functions. They are
calculated for the transformed data and then transformed back, here for
a fit of a power function:

\begin{LTXexample}[pos=t]
\begin{psgraph}[arrows=->,Dy=10](0,0)(0,0)(10,80){4.5in}{4.5in}
\psset{plotstyle=GLLSR,PowerFit,ShowEq=false}
\listplot[ShowEq,decimals=2,EqPos=1 70]{\PScatter}
\listplot[plottype=bothci,linestyle=dashed,linecolor=blue]{\PScatter}
\listplot[plottype=bothpi,linestyle=dotted,linecolor=red]{\PScatter}
\listplot[plotstyle=dots]{\PScatter}
\end{psgraph}
\end{LTXexample}

\section{Standardised residuals}
With \texttt{plottype=residuals} a cross is drawn for every data point at its
$x$ value and at the height of its standardised residual
$e_i^*=(y_i-\hat y_i)/\bigl(S_{yx}\sqrt{1-h(x_i)}\bigr)$ \cite{draper98}.
Without a systematic deviation of the data from the fit function the
residuals scatter randomly around zero and nearly all of them are between $-2$ and $2$.

\begin{LTXexample}[pos=t]
\begin{psgraph}[Dy=1,Oy=-3](0,-3)(0,-3)(10,3){4.5in}{2.5in}
\psline[linestyle=dashed](0,0)(10,0)
\psline[linestyle=dotted](0,2)(10,2)\psline[linestyle=dotted](0,-2)(10,-2)
\listplot[plotstyle=GLLSR,plottype=residuals,ShowEq=false,linecolor=red]{\Calib}
\end{psgraph}
\end{LTXexample}

\section{Plot in the transformed coordinates}
The optional argument \Lkeyword{plotcoord} selects the coordinates for the plot
and the statistics:

\medskip
\begin{tabular}{@{}lp{0.8\linewidth}@{}}\hline
\Lkeyword{plotcoord} & meaning\\\hline
\texttt{normal} & Plot in the coordinates of the data. The statistics are those
  of the transformed data (default).\\
\texttt{trans} & Plot in the transformed coordinates, where the fit is a
  straight line or a polynomial. The equation is the one of this polynomial.\\
\texttt{special} & Plot in the coordinates of the data, but with the statistics of the
  untransformed data: $\mathit{SSE}$, $S_{yy}$ and $S_{yx}$ are calculated from the values of
  $y$ itself and the intervals have the same width in $y$ direction for a given leverage.\\\hline
\end{tabular}

\medskip
The following example shows the fit of the power function from the last
but one section with \texttt{plotcoord=trans}. The axes are $\log x$ and $\log y$,
the data points are transformed with \Lcs{pstScalePoints}. This must be reset
before the fit, which transforms the data itself.

\begin{LTXexample}[pos=t]
\begin{psgraph}[arrows=->,Dx=0.2,Dy=0.5](0,0)(0,0)(1.1,2.2){4.5in}{3.5in}
\pstScalePoints(1,1){log}{log}
\listplot[plotstyle=dots]{\PScatter}
\pstScalePoints(1,1){}{}
\psset{plotstyle=GLLSR,PowerFit,plotcoord=trans,ShowEq=false}
\listplot[ShowEq,decimals=2,EqPos=0.1 2]{\PScatter}
\listplot[plottype=bothci,linestyle=dashed,linecolor=blue]{\PScatter}
\listplot[plottype=bothpi,linestyle=dotted,linecolor=red]{\PScatter}
\end{psgraph}
\end{LTXexample}

\section{Error bars with a function}
The plot style \texttt{ErrorLine} draws error bars with a length relative to
the value (\Lkeyword{relxerr}, \Lkeyword{relyerr}). With the plot style
\texttt{ErrorFunction} the length is calculated with PostScript code, given with the
optional arguments \Lkeyword{xerrfunc} and \Lkeyword{yerrfunc}. The code finds the
$x$ value (or the $y$ value) of the point on the stack and has to replace it by
the absolute error. In the example the error in $x$ is always $0.2$ and
the error in $y$ is 5\,\% of the value plus $0.3$.

\begin{LTXexample}[pos=t]
\begin{psgraph}[arrows=->,Dy=2](0,0)(0,0)(11,22){4.5in}{4.5in}
\listplot[plotstyle=ErrorFunction,arrows=|-|,arrowscale=0.5,
  xerrfunc=pop 0.2,yerrfunc=0.05 mul 0.3 add]{\Calib}
\listplot[plotstyle=dots]{\Calib}
\listplot[decimals=2,EqPos=1 20,plotstyle=GLLSR,linecolor=blue]{\Calib}
\end{psgraph}
\end{LTXexample}

\section{Numerical accuracy}
The fit is calculated by the PostScript interpreter. With \texttt{latex},
\texttt{dvips} and Ghostscript the numbers have single precision, about 7
digits; with Lua\LaTeX{} they have double precision, about 15 digits. The sums
of the least squares method need many more digits than the data has, if the
$x$ values are far from zero or if the order of the polynomial is high.
Up to version 0.04 the result with Ghostscript could then be a wrong curve,
without any warning.

Since version 0.05 the $x$ values are centred and scaled for the calculation,
$u=(x-\bar x)/s$ lies between $-1$ and $1$, and the fit is a polynomial in $u$.
The curve, the statistics and the intervals are calculated from this
polynomial. Only for the printed equation it is converted into a polynomial
in $x$. For a fit with a given \Lkeyword{Yint} the $x$ values are only scaled.
The table shows the fit of a polynomial of second order to ten points with
$x$ values from 293 to 302, with \texttt{latex}, \texttt{dvips} and Ghostscript 10.08:

\medskip
\begin{tabular}{@{}lll@{}}\hline
 & equation & $R^2$\\\hline
version 0.04 & $y=-530.786-5.02072x+0.02322x^2$ & 0.951966\\
version 0.05 & $y=68448.6-468.787x+0.802651x^2$ & 0.999716\\
exact & $y=68448.6-468.787x+0.802652x^2$ & 0.999716\\\hline
\end{tabular}

\medskip
This is always done and needs no option. It is enough for the usual fits. With
single precision it is not enough for a polynomial of a high order, which
is the case for the optional argument \Lkeyword{bignum} of the next section.

\section{Arbitrary precision arithmetic}
With the optional argument \Lkeyword{bignum} the sums, the solution of the
equations, the fitted values and the values for the intervals are calculated
with decimal numbers of the given number of digits. \texttt{bignum} without a
value means 20 digits; \texttt{bignum=0}, the default, switches it off.

\medskip
\noindent\fbox{\parbox{\dimexpr\linewidth-2\fboxsep-2\fboxrule}{%
\textbf{With Lua\LaTeX{} the option is ignored.} There the numbers have double
precision already, the fit below is correct without the option, and the
decimal numbers would be about 20 times slower. The option is meant for
\texttt{latex}, \texttt{dvips} and Ghostscript, where the numbers have single
precision. A document can use it and be compiled both ways.}}

\medskip
The show case is a polynomial of order 10 through 15 points, the values of
$y=0.8x^2/3+3\sin x+5$ for $x=1,2,\ldots,15$, rounded to one decimal. The table
shows the first coefficients of the printed equation and the standard error
of the fit, with \texttt{latex}, \texttt{dvips} and Ghostscript 10.08:

\medskip
\begin{tabular}{@{}lrrr@{}}\hline
 & without \texttt{bignum} & \texttt{bignum=20} & exact\\\hline
$a_0$ & $-3.96582$ & $-0.721988$ & $-0.721905$\\
$a_1$ & $27.269$ & $18.9622$ & $18.9620$\\
$a_2$ & $-26.0178$ & $-17.7077$ & $-17.7075$\\
$a_3$ & $14.899$ & $10.5195$ & $10.5194$\\
$a_4$ & $-5.50002$ & $-4.12873$ & $-4.12869$\\
$a_5$ & $1.28763$ & $1.0171$ & $1.01709$\\
$a_6$ & $-0.188355$ & $-0.153882$ & $-0.153881$\\
$a_7$ & $0.017063$ & $0.014233$ & $0.014233$\\
$a_8$ & $-0.000929$ & $-0.000784$ & $-0.000784$\\
$S_{yx}$ & $0.036202$ & $0.024586$ & $0.024585$\\\hline
\end{tabular}

\medskip
Without \texttt{bignum} the curve still runs close to the points, but it is
not the least squares fit: the coefficients are wrong from the first digit
and the standard error is 47\,\% too large, and with it the confidence and
prediction intervals. With \texttt{bignum=20} the coefficients are correct to
four or five digits; the rest is the precision with which the data itself
reaches the interpreter. The calculation takes longer, about one second for
this figure with Ghostscript.

\begin{LTXexample}[pos=t]
\begin{psgraph}[arrows=->,Dy=20,Dx=2](0,0)(0,0)(16,80){4.5in}{3in}
\psset{plotstyle=GLLSR,PolyOrder=10,bignum=20}
\listplot[decimals=6,EqPos=1 75,printconfig=0 1 0 0 0 0 0 1]{\HighOrder}
\listplot[plotstyle=dots]{\HighOrder}
\end{psgraph}
\end{LTXexample}

The decimal numbers come from the package \texttt{pst-numeric}, which
\texttt{pst-fit} uses if it is installed: it loads its header
\texttt{pst-numeric-bignum.pro} (not the package itself). Without
\texttt{pst-numeric}, \texttt{pst-fit} loads its own copy of the same
library, \texttt{pst-fit-bignum.pro}, so the option works either way. The
log file says which one was used. The library is \texttt{bignumlibrary.ps}
from \url{https://github.com/buddyledger/pst-numeric}; \texttt{pst-numeric}
adds matrices, complex numbers and numerical methods to it.

\subsection{The decimal numbers used directly}
The procedures of the library can also be used in own PostScript code. They
are in the dictionary \texttt{tx@NumericDict}. A decimal number is made with
\texttt{setbn} from an integer or, for all other values, from a string;
\texttt{n setbnsize} sets the number of digits; \texttt{addbn}, \texttt{subbn},
\texttt{mulbn} and \texttt{divbn} calculate, \texttt{eqbn}, \texttt{gtbn} and
\texttt{ltbn} compare, \texttt{bncvs} converts into a string and
\texttt{bncvr} into a real. There are also \texttt{pibn}, \texttt{sinbn},
\texttt{cosbn}, \texttt{tanbn}, \texttt{lnbn} and \texttt{log10bn}. The
examples use the following macro, which prints the string that the code
leaves on the stack:

\begin{lstlisting}
\newcommand\bnshow[1]{%
  \begin{pspicture}(0,-0.1)(11,0.3)
  \pscustom{\moveto(0,0)
    \code{/Courier findfont 9 scalefont setfont tx@NumericDict begin #1 end show}}
  \end{pspicture}}
\end{lstlisting}

These examples do not depend on the option \Lkeyword{bignum} and they are
calculated with Lua\LaTeX, too, where they take much longer than with
Ghostscript (the 50 digits of $\pi$ about 40 seconds instead of less than one).

\subsubsection*{50 digits of $\pi$}
For comparison the first two rows show the value which a real stores, with
single precision (Ghostscript) and with double precision (Lua\LaTeX). The
digits in gray are wrong. The last row is calculated by the library.

\begin{LTXexample}[pos=t]
\begin{tabular}{@{}ll@{}}
single precision & \texttt{\small 3.141592\wrong{7410125732421875}}\\
double precision & \texttt{\small 3.141592653589793\wrong{115997963468544}}\\
decimal, 50 digits & \bnshow{50 setbnsize pibn bncvs}
\end{tabular}
\end{LTXexample}

\subsubsection*{$\sqrt2$ with 40 digits}
Newton's method, $x\leftarrow(x+2/x)/2$, doubles the number of correct
digits with every step. Six steps from $x=1$ are enough:
\begin{LTXexample}[pos=t]
\begin{tabular}{@{}ll@{}}
single precision & \texttt{\small 1.4142135\wrong{3816986083984375}}\\
double precision & \texttt{\small 1.414213562373095\wrong{145474621858739}}\\
decimal, 40 digits & \bnshow{40 setbnsize 1 setbn
  6 { dup 2 setbn exch divbn addbn (0.5) setbn mulbn } repeat bncvs}
\end{tabular}
\end{LTXexample}

\subsubsection*{$30!$}
The factorial of 30 has 33 digits. A real keeps the first 8 or 16 of them:
\begin{LTXexample}[pos=t]
\begin{tabular}{@{}ll@{}}
single precision & \texttt{\small 26525285\wrong{1276098130314698283810816}}\\
double precision & \texttt{\small 2652528598121910\wrong{68217601719009280}}\\
decimal, 33 digits & \bnshow{33 setbnsize 1 setbn 1 1 30 { setbn mulbn } for bncvs}
\end{tabular}
\end{LTXexample}

\subsection{Pitfalls of rounding}
Every example is calculated three times. The row for single precision is the
result of Ghostscript 10.08; it cannot be calculated in this document, which
is compiled with Lua\LaTeX. The rows for double precision and for the decimal
numbers are calculated here. In every example the correct result is known,
it is given in the heading.

\subsubsection*{Ten times 0.1, minus 1, should be 0}
A real is a binary number and $0.1$ has no exact binary representation.
The error is small, but the sum is not 1 and a test for equality fails.
The decimal number $0.1$ is exact.
\begin{LTXexample}[pos=t]
\begin{tabular}{@{}ll@{}}
single precision & \texttt{\small 1.19209e-07}\\
double precision & \bnshow{0 10 { 0.1 add } repeat 1 sub 30 string cvs}\\
decimal, 20 digits & \bnshow{20 setbnsize 0 setbn 10 { (0.1) setbn addbn } repeat
  1 setbn subbn bncvs}
\end{tabular}
\end{LTXexample}

\subsubsection*{$(1+10^{-10})-1$ should be $10^{-10}$}
A small number is lost beside a large one, if the sum needs more digits than
the number has. With single precision nothing is left of $10^{-10}$.
\begin{LTXexample}[pos=t]
\begin{tabular}{@{}ll@{}}
single precision & \texttt{\small 0.0}\\
double precision & \bnshow{1 1e-10 add 1 sub 30 string cvs}\\
decimal, 20 digits & \bnshow{20 setbnsize 1 setbn (1e-10) setbn addbn 1 setbn subbn bncvs}
\end{tabular}
\end{LTXexample}

\subsubsection*{$(10^{16}+1)-10^{16}$ should be 1}
Double precision only moves the limit. This is what happens in the sums of a
least squares fit, when the $x$ values are far from zero.
\begin{LTXexample}[pos=t]
\begin{tabular}{@{}ll@{}}
single precision & \texttt{\small 0.0}\\
double precision & \bnshow{1e16 1 add 1e16 sub 20 string cvs}\\
decimal, 20 digits & \bnshow{20 setbnsize (1e16) setbn 1 setbn addbn
  (1e16) setbn subbn bncvs}
\end{tabular}
\end{LTXexample}

\subsubsection*{$(1/3)\cdot3-1$ should be 0}
Decimal numbers are rounded, too. More digits make the rounding error
smaller, but they do not remove it: $1/3$ has no exact decimal
representation. Here the reals give the correct result, because the two
rounding errors happen to cancel.
\begin{LTXexample}[pos=t]
\begin{tabular}{@{}ll@{}}
single precision & \texttt{\small 0.0}\\
double precision & \bnshow{1 3 div 3 mul 1 sub 30 string cvs}\\
decimal, 20 digits & \bnshow{20 setbnsize 1 setbn 3 setbn divbn 3 setbn mulbn
  1 setbn subbn bncvs}
\end{tabular}
\end{LTXexample}

\clearpage
\section{List of all optional arguments for \texttt{pst-fit}}

\xkvview{family=pst-fit,columns={key,type,default}}




\bgroup
\raggedright
\nocite{*}
\printbibliography
\egroup

\printindex





\end{document}

